\begin{table}[H]
\centering\footnotesize\setlength{\tabcolsep}{5pt}
\caption{Uso dos conceitos de efeito \textit{flypaper} e ilusão fiscal na literatura sobre transferências a governos locais, por periódico (2016--2026)}
\label{tab:bibliometria}
\begin{threeparttable}
\begin{tabular}{lccc}
\toprule
 & (1) & (2) & (3) \\
 & Transferências a & \textit{Flypaper} ou & (2)/(1) \\
Periódico & governos locais & ilusão fiscal & (\%) \\
\midrule
Local Government Studies & 11 & 6 & 54,5\% \\
Public Finance Review & 11 & 3 & 27,3\% \\
European Journal of Political Economy & 10 & 0 & 0,0\% \\
International Tax and Public Finance & 10 & 1 & 10,0\% \\
Public Choice & 8 & 2 & 25,0\% \\
Public Budgeting and Finance & 7 & 1 & 14,3\% \\
Lex Localis & 6 & 3 & 50,0\% \\
Regional Studies & 6 & 0 & 0,0\% \\
World Development & 6 & 0 & 0,0\% \\
Regional Science and Urban Economics & 5 & 1 & 20,0\% \\
Demais periódicos & 247 & 37 & 15,0\% \\
\midrule
Total & 327 & 54 & 16,5\% \\
\bottomrule
\end{tabular}
\begin{tablenotes}[flushleft]\footnotesize
\item Nota: A tabela resume o número de artigos publicados em periódicos científicos entre 2016 e 2026 sobre transferências intergovernamentais a governos locais. A busca foi realizada em 22/10/2026 na base Scopus, nos campos título, resumo e palavras-chave, com a expressão: ("flypaper effect" OR "fiscal illusion" OR "intergovernmental transfer*" OR "intergovernmental grant*") AND (municipal* OR "local government*" OR subnational), restrita a artigos de periódico publicados entre 2016 e 2026. A busca retornou 327 artigos. A coluna (1) apresenta todos os artigos; a coluna (2), os que tratam especificamente de efeito flypaper ou de ilusão fiscal (termos presentes no título, no resumo ou nas palavras-chave; artigos que tratam do fenômeno sem usar esses termos ficam apenas na coluna 1); a coluna (3), a participação de (2) em (1). As linhas listam os dez periódicos com mais artigos; os demais estão agregados. Os dados de 2026 são parciais (artigos indexados até a data da busca). Processamento: pacote bibliometrix (Aria e Cuccurullo, 2017) no R.
\end{tablenotes}
\end{threeparttable}
\end{table}
